\begingroup
\section{Node ZG: strong ZG-T1 retracted; structured ZG-COM closed}
\label{module:ZG}
\noindent\textit{Audited source: \nolinkurl{ZG_zero_gap.tex}.}
\subsection{Intrinsic cut and source decomposition}

Let $I_i=[\theta_i,\theta_{i+1}]$ have length $h_i>0$ and midpoint
$c_i$.  Put
\begin{equation}\label{ZG:eq:q}
 q_h=\sum_iq_i,
 \qquad
 q_i(\theta_i+x)=\frac12\min\{x^2,(h_i-x)^2\}
 \quad(0\le x\le h_i),
\end{equation}
where every $q_i$ is extended by zero outside $I_i$.  Define
\begin{equation}\label{ZG:eq:T}
 T(h)=\sum_i h_i|\Hil q_h(c_i)|^2.
\end{equation}

As in stage 153, let
\begin{equation}\label{ZG:eq:D}
 \ell^2=\frac1{2\pi}\sum_i h_i^3,
 \qquad
 \Dvar(h)=\sum_i h_i(h_i^2-\ell^2)^2.
\end{equation}
The far region $\Dvar\ge\delta\sum_i h_i^5$ is already closed by the
absolute adaptive sampling estimate.  In the complementary region choose
$0<\rho<1/8$ and set
\begin{equation}\label{ZG:eq:GB}
 \mathcal G=\{i:|h_i/\ell-1|<\rho\},
 \qquad \mathcal B=\{1,\ldots,N\}\setminus\mathcal G.
\end{equation}
Put
\begin{equation}\label{ZG:eq:bE}
 M=|\mathcal G|,\qquad b=\frac{2\pi}{M},\qquad
 E=\sum_{i\in\mathcal B}h_i.
\end{equation}
The near-stratum selection gives $b\asymp_\rho\ell$ and $E$ smaller than
a fixed multiple of $2\pi$.  We fix the stratum threshold so that
$E\le\pi/4$.  If $E$ is larger, the absolute sampling bound
$T\le C\sum_i h_i^5$, together with \eqref{ZG:eq:budget} below and
$b^4\le Cb^4E$, already belongs to the closed far row.  The following
static budget is the only
place where the definition of the bad set is used:
\begin{equation}\label{ZG:eq:budget}
 b^3\sum_{i\in\mathcal G}(h_i-b)^2
 +b^4E+\sum_{i\in\mathcal B}h_i^5
 \le C_\rho\Dvar(h).
\end{equation}
Indeed, the good term is the reference-charge estimate of stage 153.
On the bad set, with $r=h_i/\ell$,
\[
 |r-1|\ge\rho,
 \qquad r^5+r\le C_\rho r(r^2-1)^2.
\]
Multiplication by $\ell^5$, summation, and $b\asymp_\rho\ell$ give
\[
 b^4E+\sum_{i\in\mathcal B}h_i^5
 \le C_\rho\sum_{i\in\mathcal B}
 h_i(h_i^2-\ell^2)^2,
\]
which proves \eqref{ZG:eq:budget} without invoking the retracted block
lemma of stage 155.

Split
\begin{equation}\label{ZG:eq:split}
 q_{\rm G}=\sum_{i\in\mathcal G}q_i,
 \qquad q_{\rm B}=\sum_{i\in\mathcal B}q_i.
\end{equation}
Then
\begin{equation}\label{ZG:eq:Tsplit}
 T(h)\le2T_{\rm G}+2T_{\rm B},
 \qquad
 T_X:=\sum_i h_i|\Hil q_X(c_i)|^2.
\end{equation}

\subsection{The bad source is an absolute square}

We recall the special adaptive sampling fact proved in stage 153.  If an
interval partition has midpoints $c_i$ and $f\in L^2(\T)$ satisfies
\begin{equation}\label{ZG:eq:selfzero}
 \operatorname {pv}\int_{I_i}\frac1{2\pi}
 \cot\frac{c_i-y}{2}f(y)\,dy=0
 \quad\text{for every }i,
\end{equation}
then
\begin{equation}\label{ZG:eq:adaptive}
 \sum_i h_i|\Hil f(c_i)|^2\le C\|f\|_2^2.
\end{equation}
It follows by comparing the truncation at radius $h_i/2$ with the
maximal Hilbert transform and integrating the Cotlar inequality over the
disjoint cells.

\begin{lemma}[Bad-source square]\label{ZG:lem:bad}
For the decomposition \eqref{ZG:eq:split},
\begin{equation}\label{ZG:eq:bad}
 \boxed{
 T_{\rm B}\le C\sum_{j\in\mathcal B}h_j^5.}
\end{equation}
\end{lemma}

\begin{proof}
If $i\in\mathcal G$, the function $q_{\rm B}$ vanishes on $I_i$.  If
$i\in\mathcal B$, its restriction to $I_i$ is the even bubble $q_i$,
while the cotangent kernel centred at $c_i$ is odd.  Thus
\eqref{ZG:eq:selfzero} holds for $f=q_{\rm B}$ on every target cell.
Equation \eqref{ZG:eq:adaptive} and the exact cell integral
\[
 \|q_{\rm B}\|_2^2
 =\sum_{j\in\mathcal B}\int_{I_j}q_j^2
 =\frac1{320}\sum_{j\in\mathcal B}h_j^5
\]
prove \eqref{ZG:eq:bad}.
\end{proof}

\begin{remark}
This single application of \eqref{ZG:eq:adaptive} is the decisive cleaning
step.  All bad midpoints are now inside the source $q_{\rm B}$ and are
paid by their individual $h_j^5$ energies.  A maximal bad block is never
asked to pack its internal midpoint set at every dyadic scale.
\end{remark}

\subsection{The good skeleton with zero-source gaps}

Delete the bad bubbles but retain their intervals as gaps.  Thus
$q_{\rm G}$ consists of comparable good bubbles, while a maximal
consecutive bad block $B_\nu$ is a zero-source gap of length
\begin{equation}\label{ZG:eq:gaps}
 e_\nu=\sum_{i\in B_\nu}h_i,
 \qquad \sum_\nu e_\nu=E.
\end{equation}
Only the two endpoints $\partial B_\nu$ are singular locations for the
good source.

For a target interval $Q$ of physical length $L_Q$ (either a member of
one of the shifted grids or a node of the balanced target tree below),
define the endpoint weight
\begin{equation}\label{ZG:eq:Omega-end}
 \Omega_{\rm end}(B,Q)=
 \begin{cases}
 \displaystyle \frac e{L_Q}
 \left(1+\frac{\operatorname {dist}(B,Q)}{L_Q}\right)^{-3},
 &e\le L_Q,\\[3mm]
 \displaystyle \frac {L_Q}e
 \left(1+\frac{\operatorname {dist}(\partial B,Q)}{L_Q}\right)^{-3},
 &e>L_Q.
 \end{cases}
\end{equation}

\begin{lemma}[Genuine endpoint packing]\label{ZG:lem:end-packing}
For disjoint gaps, on each shifted dyadic grid and on every balanced
target interval tree used below,
\begin{equation}\label{ZG:eq:end-packing}
 \sup_Q\sum_\nu\Omega_{\rm end}(B_\nu,Q)
 +\sup_\nu\sum_Q\Omega_{\rm end}(B_\nu,Q)\le C.
\end{equation}
\end{lemma}

\begin{proof}
For $e\le L_Q$, disjointness bounds the total gap length in the $n$th
distance annulus by $C2^nL_Q$; the factor
$L_Q^{-1}2^{-3n}$ is summable.  For $e>L_Q$, sort gaps by
$e\asymp2^mL_Q$.  The endpoints themselves need not be
$2^mL_Q$-separated: two large gaps can be separated by one good cell.
What is separated is the interior mass of the disjoint gaps.  Hence an
annulus of length $C2^nL_Q$ meets at most
$C(1+2^{n-m})$ such gaps.  Multiplication by
$2^{-m}2^{-3n}$ is summable in $m,n$.

With a gap fixed, the coarse scale sum is
$\sum_{L_Q\ge e}e/L_Q<\infty$.  At fine scales there are only two
singular points, so the annular sum is bounded at each scale and
$\sum_{L_Q<e}L_Q/e<\infty$.  This proves \eqref{ZG:eq:end-packing}.  The
argument would fail with all internal bad midpoints in place of the two
endpoints; Lemma~\ref{ZG:lem:bad} is what removed them.
\end{proof}

We next isolate the precise zero-gap extension of the variable-grid
square.  Only the good midpoint atoms occur in the opening operator:
\begin{equation}\label{ZG:eq:tau}
 d\tau_{\rm G}=\sum_{i\in\mathcal G}h_i\delta_{c_i},
 \qquad d\mu=dx.
\end{equation}
Thus the source and target measures are different.  In particular the
one-weight matrix theorem of stage 154 cannot be quoted here.  We give
the continuous-to-atomic operator and its square estimate explicitly.

The exact identity behind the next estimate is worth recording.  With
$\Phi(x)=\sum_{n\ge1}\sin(nx)/n^2$, twice integrating
$q_j''=\mathbf1_{I_j}-h_j\delta_{c_j}$ gives
\begin{equation}\label{ZG:eq:Peano-H}
 \Hil q_{\rm G}(x)=\frac1\pi\sum_{j\in\mathcal G}
 \left\{h_j\Phi(x-c_j)-\int_{I_j}\Phi(x-y)\,dy\right\}.
\end{equation}
Each bracket has zero mass and zero first moment.  Formula
\eqref{ZG:eq:Peano-H}, rather than separate endpoint kernels, is used until
the first source--target separating scale.

Put $B=\bigcup_\nu B_\nu$, $G^*=\bigcup_{j\in\mathcal G}I_j$, and
\begin{equation}\label{ZG:eq:sigma}
 \gamma=\frac EM,
 \qquad
 \sigma=\mathbf1_B-
 \sum_{j\in\mathcal G}\frac{\gamma}{h_j}\mathbf1_{I_j}.
\end{equation}
Then $\int_\T\sigma=0$.  Let $u$ be the periodic primitive of $\sigma$
with zero mean and set
\begin{equation}\label{ZG:eq:collapse-map}
 R_t(x)=x-tu(x),\qquad 0\le t\le1.
\end{equation}
On a good cell $R_t'=1+t\gamma/h_j$; on a gap $R_t'=1-t$.
Consequently $R_1$ collapses every gap and maps $I_j$ affinely onto an
interval of length $h_j+\gamma$.  In particular
$R_1(c_j)=\widehat c_j$.

For $f\in L^2(\mu)$ put
\[
 f_0=f-\frac1{2\pi}\int_\T f.
\]
For an oriented cyclic arc from $x$ to $z$, let
$\epsilon_{x,z}$ be its signed indicator.  If $V_f$ is the periodic
zero-mean primitive of $f_0$, then
\begin{equation}\label{ZG:eq:oriented-primitive}
 V_f(z)-V_f(x)=\int_\T\epsilon_{x,z}(s)f_0(s)\,ds.
\end{equation}
For $y\in I_j$, $j\in\mathcal G$, write $m(y)=c_j$ and
$a_t(y)=R_t'(y)$.  Define
\begin{align}
 P_t(x,y)&=\Phi(R_t x-R_t m(y))-\Phi(R_t x-R_t y),
 \label{ZG:eq:Pt}\\
 (\mathcal T_tf)(x)
 &=\frac1\pi\int_{G^*}
 \Big[-f_0(y)P_t(x,y)\notag\\
 &\hspace{13mm}+a_t(y)\{\Phi'(R_tx-R_tm(y))
       [V_f(m(y))-V_f(x)]\notag\\
 &\hspace{35mm}-\Phi'(R_tx-R_ty)
       [V_f(y)-V_f(x)]\}\Big],dy .
 \label{ZG:eq:Tt}
\end{align}
This is a linear map from $L^2(\mu)$ to functions on the good target
points and annihilates constants.  On the zero-mean subspace,
\eqref{ZG:eq:oriented-primitive} gives the raw kernel
\begin{align}
 K_t(x,s)=\frac1\pi\Big[&-\mathbf1_{G^*}(s)P_t(x,s)
 \notag\\
 &+\int_{G^*}a_t(y)
 \{\Phi'(R_tx-R_tm(y))\epsilon_{x,m(y)}(s)\notag\\
 &\hspace{35mm}-\Phi'(R_tx-R_ty)\epsilon_{x,y}(s)\}\,dy\Big].
 \label{ZG:eq:Kt}
\end{align}
For arbitrary $f$ the actual kernel is the source-centred version
\begin{equation}\label{ZG:eq:Kt-centred}
 K_t^0(x,s)=K_t(x,s)-\frac1{2\pi}\int_\T K_t(x,r)\,dr,
 \qquad
 \mathcal T_tf(x)=\int_\T K_t^0(x,s)f(s)\,ds.
\end{equation}
Below we suppress the superscript $0$.  The subtracted row constant has
zero source difference.  Integrating \eqref{ZG:eq:Kt} in $s$, using
$\int_\T\Phi=0$, and pairing the two terminal half-cells gives
\begin{equation}\label{ZG:eq:row-gauge}
 \left|\frac1{2\pi}\int_\T K_t(x,s)\,ds\right|
 +\left|\partial_x\frac1{2\pi}\int_\T K_t(x,s)\,ds\right|
 \le C_\rho b^2.
\end{equation}
Thus the centred kernel obeys the same size, difference, and testing
bounds (with a changed constant).

\begin{proposition}[The strong ZG--$T1$ square is false]
\label{ZG:prop:Peano-counterexample}
There is no constant independent of $M$ for which
\begin{equation}\label{ZG:eq:false-Peano-square}
 \sum_{i\in\mathcal G}h_i|\mathcal T_0f(c_i)|^2
 \le C b^4\|f\|_{L^2(\T)}^2
\end{equation}
holds for every $f$.  In fact, on the uniform gap-free fan the norm of
$\mathcal T_0:L^2(dx)\to L^2(\tau_{\rm G})$ is at least a constant
multiple of $b$.
\end{proposition}

\begin{proof}
Take $h_j=b=2\pi/M$, $c_j=(j+\tfrac12)b$, and
\[
 f_M(\theta)=(2\pi)^{-1/2}e^{i(M-1)\theta}.
\]
Put $n=M-1$ and $q=e^{-ib}$.  Cellwise integration by parts in
\eqref{ZG:eq:Tt}, followed by cyclic summation, gives for the unnormalised
mode $e^{in\theta}$
\begin{equation}\label{ZG:eq:counterexample-exact}
 \mathcal T_0(e^{in\cdot})(c_0)
 =\frac{B_M+C_M}{\pi i(M-1)},
\end{equation}
where
\begin{align}
 B_M&=-(1-q)\sum_{j=0}^{M-1}q^j\Phi(jb),
 \label{ZG:eq:counterexample-B}\\
 C_M&=b e^{-ib/2}\sum_{j=1}^{M-1}
       \Phi'(jb)(1-q^j).
 \label{ZG:eq:counterexample-C}
\end{align}
Here is the calculation before summing.  Write $x=c_0$, $m=c_j$,
$d=x-m=-jb$, $F(y)=e^{iny}$, and
$P(r)=\Phi(d)-\Phi(d-r)$, and denote by
$[\mathcal T_0F(x)]_{I_j}$ the $I_j$-integral in \eqref{ZG:eq:Tt}.
Integration by parts of $-FP$, using
$F=(F/in)'$, cancels its interior $F\Phi'$ term against the last
primitive term in \eqref{ZG:eq:Tt} and gives
\begin{multline}\label{ZG:eq:counterexample-cell}
 \pi[\mathcal T_0F(x)]_{I_j}
 =\frac1{in}\Big\{
 -[F(m+r)P(r)]_{-b/2}^{b/2}\\
 +F(x)[\Phi(d+b/2)-\Phi(d-b/2)]
 +b\Phi'(d)[F(m)-F(x)]\Big\}.
\end{multline}
For $j=0$ the last product is understood by continuity and is zero.
Now
\[
 F(c_j)=-e^{-ib/2}q^j,
 \qquad \Phi(-z)=-\Phi(z),
 \qquad \Phi'(-z)=\Phi'(z).
\]
Substitution in \eqref{ZG:eq:counterexample-cell}, followed by one cyclic
index shift in the endpoint terms, gives
\eqref{ZG:eq:counterexample-B}--\eqref{ZG:eq:counterexample-C} and hence
\eqref{ZG:eq:counterexample-exact}.

The Riemann sum in \eqref{ZG:eq:counterexample-B} satisfies
\[
 b\sum_{j=0}^{M-1}e^{-ijb}\Phi(jb)
 \longrightarrow\int_0^{2\pi}e^{-ix}\Phi(x)\,dx=-i\pi.
\]
Since $1-e^{-ib}=ib+O(b^2)$, it follows that $B_M\to-\pi$.
Although $\Phi'$ has logarithmic endpoint singularities, the function
\[
 x\longmapsto\Phi'(x)(1-e^{-ix})
\]
extends continuously by zero at both endpoints.  Hence
\begin{align*}
 C_M&\longrightarrow
 \int_0^{2\pi}\Phi'(x)(1-e^{-ix})\,dx\\
 &= -\int_0^{2\pi}\Phi'(x)e^{-ix}\,dx=-\pi;
\end{align*}
the last equality is the first cosine coefficient of
$\Phi'(x)=\sum_{k\ge1}\cos(kx)/k$.  Thus
\begin{equation}\label{ZG:eq:counterexample-asymptotic}
 |\mathcal T_0(e^{in\cdot})(c_0)|\sim\frac2M.
\end{equation}
Translation covariance gives the same modulus at all $M$ target atoms.
The normalisation of $f_M$ cancels the total target mass $2\pi$, and so
\[
 \|\mathcal T_0f_M\|_{L^2(\tau_{\rm G})}
 =\frac{|B_M+C_M|}{\pi(M-1)}\sim\frac2M.
\]
Since $b^2=4\pi^2/M^2$, the ratio of the two sides at the norm level is
asymptotic to $M/(2\pi^2)$.  This contradicts
\eqref{ZG:eq:false-Peano-square}.  If a nonempty deleted block is required,
insert one gap of length $o(M^{-3})$ and distribute that length over the
good cells.  For each fixed $M$ the finite cell integrals above depend
continuously on those endpoints, so the same contradiction persists.
\end{proof}

\begin{remark}[Consequence for the former proof route]
The former ZG--$T1$ obligation asserted \eqref{ZG:eq:false-Peano-square} for
all $f$ and all $t$.  Proposition~\ref{ZG:prop:Peano-counterexample} shows
that no completion of its proposed Haar paragraph can prove that claim.
In particular, the previously stated kernel/testing package cannot have
all the asserted uniform consequences.  None of those claims is used
below.  The opening argument only needs the special density $\sigma$ and
the time integral appearing in the exact deformation identity.
\end{remark}

\begin{lemma}[Endpoint Calder\'on estimate for a small change of variable]
\label{ZG:lem:small-change}
Let $S:\T\to\T$ be an orientation-preserving, degree-one, piecewise
affine map, and choose its lift in the form $S(x)=x+v(x)$.  Suppose
\begin{equation}\label{ZG:eq:small-change-hypothesis}
 \|v'\|_\infty\le\eta<1.
\end{equation}
For $f\in W^{1,\infty}(\T)$ put
\[
 \mathfrak C_Sf=(\Hil f)\circ S-\Hil(f\circ S).
\]
Then
\begin{align}
 \|\mathfrak C_Sf\|_2
 &\le C_\eta\|S'-1\|_2\|f\|_\infty,
 \label{ZG:eq:small-change-H0}\\
 \|(\mathfrak C_Sf)'\|_2
 &\le C_\eta\|S'-1\|_2\|f'\|_\infty.
 \label{ZG:eq:small-change-H1}
\end{align}
The same constants hold for limits of such maps with nonnegative
derivative, provided each approximating map satisfies
\eqref{ZG:eq:small-change-hypothesis}.
\end{lemma}

\begin{proof}
We include the endpoint argument because replacing the right side of
\eqref{ZG:eq:small-change-H0} by an arbitrary-input $L^2$ norm is precisely
the false step retracted above.  First smooth $v$ and $f$; the estimates
obtained below are uniform under this smoothing.

On a lifted interval, changing variables in the first Hilbert transform
and writing
\[
 \Delta_v(x,y)=\frac{v(x)-v(y)}{x-y}
\]
gives, for the singular part of the cotangent kernel,
\begin{equation}\label{ZG:eq:calderon-kernel}
 \frac{1+v'(y)}{Sx-Sy}-\frac1{x-y}
 =\frac{v'(y)-\Delta_v(x,y)}
 {(x-y)(1+\Delta_v(x,y))}.
\end{equation}
For the moment suppose $\|v'\|_\infty\le\eta_0$, where the universal
$\eta_0>0$ will be chosen below.  Expand the last denominator
geometrically.  It is therefore enough to estimate
\begin{equation}\label{ZG:eq:calderon-Am}
 A_m(v,g)(x)=\operatorname {pv}\!\int
 g(y)\frac{(v'(y)-\Delta_v(x,y))
                 \Delta_v(x,y)^m}{x-y}\,dy,
 \qquad m\ge0,
\end{equation}
where in the application $g=f\circ S$.  The following two endpoint
multiplier bounds are the elementary Calder\'on estimate:
\begin{align}
 \|A_m(v,g)\|_2
 &\le C^{m+1}\|v'\|_2\|v'\|_\infty^m\|g\|_\infty,
 \label{ZG:eq:Am-H0}\\
 \|\partial_xA_m(v,g)\|_2
 &\le C^{m+1}\|v'\|_2\|v'\|_\infty^m\|g'\|_\infty.
 \label{ZG:eq:Am-H1}
\end{align}
Here is a direct frequency verification.  Use
\[
 \Delta_v(x,y)=\int_0^1v'(y+s(x-y))\,ds
\]
in \eqref{ZG:eq:calderon-Am}, expand into Fourier series, and divide the
frequencies into the regions in which one of the $m+1$ copies of $v'$
has largest modulus.  The difference
$v'(y)-\Delta_v(x,y)$ makes the multiplier vanish unless such a
$v'$-frequency dominates the frequency of $g$.  On each region the
multiplier and all its discrete differences are bounded by $C^{m+1}$.
Smooth dyadic resolution consequently writes that piece as
\[
 \sum_r P_rv'\,
 T_r(P_{\le r+3}v',\ldots,P_{\le r+3}v',P_{\le r+3}g),
\]
up to a uniformly summable family of translates; the kernels of $T_r$
have uniformly bounded $L^1$ norm.  The square-function inequality for
the sole high-frequency factor and the uniform $L^\infty$ bounds for
smooth low-frequency projections prove \eqref{ZG:eq:Am-H0}.  After applying
$\partial_x$, summation by parts in the Fourier frequency of $g$ moves
the output frequency onto $g'$; on the nonzero multiplier support its
coefficient is bounded by the same largest $v'$-frequency.  The identical
decomposition proves \eqref{ZG:eq:Am-H1}.  This also shows that the constants
are independent of the number of affine pieces.

Choose $\eta_0$ so that $C\eta_0<1/2$.  Summing
\eqref{ZG:eq:Am-H0}--\eqref{ZG:eq:Am-H1} with the geometric coefficients gives
the assertions for the flat singular kernel whenever
$\|v'\|_\infty\le\eta_0$.  The
difference $\frac12\cot(t/2)-1/t$ is smooth.  Subtracting $g(x)$ (constants
are annihilated) and integrating once by parts reduces its two estimates
to convolution with an $L^1$ kernel, and gives the same right sides.
Thus \eqref{ZG:eq:small-change-H0}--\eqref{ZG:eq:small-change-H1} hold in the
small-distortion case.  Notice that
\[
 \|g\|_\infty=\|f\|_\infty,
 \qquad \|g'\|_\infty\le(1+\eta_0)\|f'\|_\infty.
\]

It remains to pass from $\eta_0$ to the stated fixed $\eta<1$ without
using the divergent series above.  For $0\le s\le1$ define the degree-one
map $S_s$ by
\begin{equation}\label{ZG:eq:logarithmic-interpolation}
 S_s'(x)=\frac{S'(x)^s}{(2\pi)^{-1}\int_\T S'(y)^s\,dy},
 \qquad S_0=\operatorname {id}.
\end{equation}
Up to a harmless rotation, $S_1=S$.  Partition $[0,1]$ into a number
$L=L(\eta,\eta_0)$ of equal pieces and set
$T_j=S_{(j+1)/L}\circ S_{j/L}^{-1}$.  Since
$1-\eta\le S'\le1+\eta$, differentiating
\eqref{ZG:eq:logarithmic-interpolation} in $s$ gives
\begin{equation}\label{ZG:eq:interpolation-cost}
 \|T_j'-1\|_\infty\le\eta_0,
 \qquad
 \sum_{j=0}^{L-1}\|T_j'-1\|_2
 \le C_\eta\|S'-1\|_2.
\end{equation}
Indeed, on this compact interval the functions $\log t$ and $t-1$ are
uniformly comparable, and change of variables by every $S_s$ has
Jacobian bounded above and below by constants depending only on $\eta$.

For degree-one maps the commutator has the exact cocycle
\begin{equation}\label{ZG:eq:change-cocycle}
 \mathfrak C_{S_2\circ S_1}f
 =\mathfrak C_{S_1}(f\circ S_2)
   +(\mathfrak C_{S_2}f)\circ S_1.
\end{equation}
Iterate \eqref{ZG:eq:change-cocycle} over the $T_j$.  All intermediate
Jacobians and inverse Jacobians are bounded by $C_\eta$; composition
preserves the $L^\infty$ norm and multiplies the derivative norm by at
most $C_\eta$.  Apply the already proved small-distortion estimates to
each summand and use \eqref{ZG:eq:interpolation-cost}.  This proves
\eqref{ZG:eq:small-change-H0}--\eqref{ZG:eq:small-change-H1} for every fixed
$\eta<1$.

Approximation and weak lower semicontinuity in $H^1$ remove the smoothing
and prove the lemma.
\end{proof}

Put $R=R_1$ and $\widehat h_j=h_j+\gamma$ on the good indices.  After
deleting the gaps, the $\widehat h_j$ form a positive partition of the
circle and remain comparable with $b$.  Let $\widehat q$ be its complete
quadratic spline.  Set $p=\widehat q\circ R$ and define the monotone-pullback
commutator
\begin{equation}\label{ZG:eq:collapse-commutator}
 \mathfrak C_R\widehat q(x)
 :=\Hil\widehat q(Rx)-\Hil p(x).
\end{equation}
Changing variables $z=R(y)$ in the first Hilbert transform gives the
entirely explicit kernel formula
\begin{multline}\label{ZG:eq:collapse-commutator-kernel}
 \mathfrak C_R\widehat q(x)=\frac1{2\pi}\int_\T
 \widehat q(Ry)\Big\{R'(y)
 \cot\frac{Rx-Ry}{2}\\
 -\cot\frac{x-y}{2}\Big\}\,dy ,
\end{multline}
with principal values taken symmetrically.  Here $R'=0$ on every gap and
$R'=1+\gamma/h_j$ on $I_j$.

\begin{theorem}[Collapsed-bubble commutator square]
\label{ZG:thm:collapse-commutator}
One has
\begin{equation}\label{ZG:eq:collapse-commutator-square}
 \boxed{
 \sum_{i\in\mathcal G}h_i
 |\mathfrak C_R\widehat q(c_i)|^2
 \le C_\rho b^4E.}
\end{equation}
The constant must be independent of the number, length, and subdivision
of the gaps.
\end{theorem}

\begin{proof}
We resolve the collapse into geometrically summable, nondegenerate
changes of variable.  If $B_\nu$ is a gap of length $e_\nu$, define the
level-$k$ geometry by
\begin{equation}\label{ZG:eq:dyadic-geometry}
 e_\nu^{(k)}=2^{-k}e_\nu,
 \qquad h_j^{(k)}=h_j+(1-2^{-k})\gamma,
 \qquad k\ge0.
\end{equation}
These lengths sum to $2\pi$.  Let $S_k$ be the affine map from the
level-$k$ circle to the level-$(k+1)$ circle which preserves the cyclic
order of all cells.  Thus
\begin{equation}\label{ZG:eq:Sk-derivative}
 S_k'=\frac12\quad\hbox{on the gaps},
 \qquad
 S_k'=1+\frac{2^{-k-1}\gamma}{h_j^{(k)}}
       \quad\hbox{on the $j$th good cell}.
\end{equation}
Put $R_0=\operatorname {id}$ and $R_{k+1}=S_k\circ R_k$.  Then
$R_k$ converges uniformly to the collapse $R$.

The good lengths are mutually comparable and their average is
$g=(2\pi-E)/M$.  Since $E\le\pi/4$ and $\rho<1/8$, it follows directly
from \eqref{ZG:eq:Sk-derivative} that
\begin{equation}\label{ZG:eq:Sk-uniform}
 \|S_k'-1\|_\infty\le\frac12,
 \qquad
 1\le R_k'|_{I_j}\le C_\rho .
\end{equation}
Moreover, using $\sum_j(h_j^{(k)})^{-1}\le C_\rho M/b$ and
$\gamma=E/M$, we have the exact layer computation
\begin{align}
 \|S_k'-1\|_2^2
 &=\frac14,2^{-k}E
   +2^{-2k-2}\gamma^2
       \sum_{j\in\mathcal G}\frac1{h_j^{(k)}}
 \le C_\rho2^{-k}E.                         \label{ZG:eq:Sk-L2}
\end{align}

Let $R_{k,\infty}$ be the monotone tail collapse from the level-$k$
geometry to the final complete fan, and put
\begin{equation}\label{ZG:eq:pk}
 p_k=\widehat q\circ R_{k,\infty}.
\end{equation}
Then $p_0=p$, $p_k=p_{k+1}\circ S_k$, and $p_k$ vanishes on the
level-$k$ gaps.  On a good cell $R_{k,\infty}$ is affine with derivative
$\widehat h_j/h_j^{(k)}$.  Since all these lengths are comparable with
$b$, the explicit quadratic bubble gives, uniformly in $k$,
\begin{equation}\label{ZG:eq:pk-bounds}
 \|p_k\|_\infty\le C_\rho b^2,
 \qquad \|p_k'\|_\infty\le C_\rho b.
\end{equation}

On the original good set define
\[
 G_k=(\Hil p_k)\circ R_k.
\]
The identities $R_{k+1}=S_k\circ R_k$ and
$p_k=p_{k+1}\circ S_k$ give the exact commutator cocycle
\begin{equation}\label{ZG:eq:commutator-cocycle}
 G_{k+1}-G_k=(\mathfrak C_{S_k}p_{k+1})\circ R_k.
\end{equation}
Lemma~\ref{ZG:lem:small-change}, \eqref{ZG:eq:Sk-L2},
\eqref{ZG:eq:pk-bounds}, and the change of variables on the good cells in
\eqref{ZG:eq:Sk-uniform} therefore imply
\begin{align}
 \|G_{k+1}-G_k\|_{L^2(G)}
 &\le C_\rho b^2 2^{-k/2}E^{1/2},
 \label{ZG:eq:Gk-H0}\\
 \|(G_{k+1}-G_k)'\|_{L^2(G)}
 &\le C_\rho b 2^{-k/2}E^{1/2}.
 \label{ZG:eq:Gk-H1}
\end{align}
Both series are summable.  Since $G_0=\Hil p$ and
$G_k\to(\Hil\widehat q)\circ R$ in $H^1$ on the union of the original
good cells, their sum yields
\begin{equation}\label{ZG:eq:commutator-continuous-H1}
 \|\mathfrak C_R\widehat q\|_{L^2(G)}^2
 +b^2\|(\mathfrak C_R\widehat q)'\|_{L^2(G)}^2
 \le C_\rho b^4E.
\end{equation}

For $F\in H^1(I_i)$,
\[
 h_i|F(c_i)|^2
 \le2\int_{I_i}|F|^2+h_i^2\int_{I_i}|F'|^2.
\]
Apply this to $F=\mathfrak C_R\widehat q$, sum over the disjoint good
cells, and use $h_i\asymp_\rho b$ together with
\eqref{ZG:eq:commutator-continuous-H1}.  This proves
\eqref{ZG:eq:collapse-commutator-square}.
\end{proof}

\begin{lemma}[Static zero-gap opening square]\label{ZG:lem:opening}
Let
$\widehat c_j$ be the midpoints of the preceding collapsed partition and put
\begin{equation}\label{ZG:eq:reference-vector}
 U_i=\Hil\widehat q(\widehat c_i),\qquad i\in\mathcal G.
\end{equation}
Then
\begin{equation}\label{ZG:eq:opening-square}
 \boxed{
 \sum_{i\in\mathcal G}h_i
 \left|\Hil q_{\rm G}(c_i)-U_i\right|^2
 \le C_\rho b^4E.}
\end{equation}
The constant is independent of the number and subdivision of the gaps.
\end{lemma}

\begin{proof}
For $x=c_i$ define
\begin{equation}\label{ZG:eq:Wt}
 W_i(t)=\frac1\pi\sum_{j\in\mathcal G}\int_{I_j}a_t(y)
 \{\Phi(R_tc_i-R_tc_j)-\Phi(R_tc_i-R_ty)\}\,dy.
\end{equation}
Changing variables $z=R_t(y)$ shows that $W_i(t)$ is the Hilbert
sample of the quadratic bubbles on the images $R_t(I_j)$, evaluated at
their transported midpoint.  Hence
\[
 W_i(0)=\Hil q_{\rm G}(c_i),
 \qquad W_i(1)=\Hil\widehat q(\widehat c_i)=U_i.
\]
Differentiating \eqref{ZG:eq:Wt} under the integral and using $u'=\sigma$
gives the exact identity
\begin{equation}\label{ZG:eq:opening-identity}
 \dot W_i(t)=\mathcal T_t\sigma(c_i).
\end{equation}
There is no omitted endpoint term: the gap motion is contained in the
primitive differences in \eqref{ZG:eq:Tt}.

It remains to separate the part of the opening error already controlled
by elementary Sobolev sampling.  On a good cell put
$a_j=1+\gamma/h_j=\widehat h_j/h_j$.  The quadratic bubble scales
exactly, so
\begin{equation}\label{ZG:eq:pullback-bubble}
 p|_{I_j}=a_j^2q_j,
 \qquad p=0\quad\hbox{on every gap}.
\end{equation}
Consequently $r:=q_{\rm G}-p$ belongs to $H^1(\T)$ and, since
$h_j\asymp_\rho b$ and the near stratum has bounded small $E$,
\begin{align}
 \|r\|_2^2
 &\le C_\rho E^2\sum_{j\in\mathcal G}h_j^5
 \le C_\rho b^4E,\label{ZG:eq:amplitude-H0}\\
 \|r'\|_2^2
 &\le C_\rho E^2\sum_{j\in\mathcal G}h_j^3
 \le C_\rho b^2E.\label{ZG:eq:amplitude-H1}
\end{align}
Indeed $|a_j^2-1|\le C_\rho E$, while
$Mb=2\pi$.  For every $F\in H^1(I_j)$ the fundamental theorem of
calculus and Cauchy--Schwarz give
\[
 h_j|F(c_j)|^2
 \le2\int_{I_j}|F|^2+h_j^2\int_{I_j}|F'|^2.
\]
Apply this with $F=\Hil r$, sum the good cells, and use the $L^2$
isometry of $\Hil$ and its commutation with differentiation.  Equations
\eqref{ZG:eq:amplitude-H0}--\eqref{ZG:eq:amplitude-H1} yield
\begin{equation}\label{ZG:eq:amplitude-trace}
 \sum_{i\in\mathcal G}h_i|\Hil r(c_i)|^2
 \le C_\rho b^4E.
\end{equation}
Finally, \eqref{ZG:eq:collapse-commutator} gives the exact decomposition
\[
 \Hil q_{\rm G}(c_i)-\Hil\widehat q(Rc_i)
 =\Hil r(c_i)-\mathfrak C_R\widehat q(c_i).
\]
Combine \eqref{ZG:eq:amplitude-trace} with
\eqref{ZG:eq:collapse-commutator-square} to prove
\eqref{ZG:eq:opening-square}.  Integrating
\eqref{ZG:eq:opening-identity} also shows that this static difference is
$-\int_0^1\mathcal T_t\sigma(c_i)\,dt$; hence ZG--COM closes exactly the
structured integrated opening-flow square, without a pointwise-in-$t$
operator estimate.

\end{proof}

The restriction $i\in\mathcal G$ in Lemma~\ref{ZG:lem:opening} is
essential.  On a bad midpoint the collapsed reference value is an
endpoint trace, not zero.  It must be retained in the endpoint layer,
rather than incorrectly charged to the opening density alone.

\begin{lemma}[Normalized-bubble bad-target trace]\label{ZG:lem:gap-target}
With the same notation,
\begin{equation}\label{ZG:eq:gap-target}
 \boxed{
 \sum_{i\in\mathcal B}h_i|\Hil q_{\rm G}(c_i)|^2
 \le C_\rho\left\{
 b^3\sum_{j\in\mathcal G}(\widehat h_j-b)^2+b^4E\right\}.}
\end{equation}
The constant is independent of the number and subdivision of the gaps.
\end{lemma}

\begin{proof}
Put
\[
 g:=\frac{2\pi-E}{M}=b-\gamma.
\]
Thus $\widehat h_j-b=h_j-g$, and $g\asymp_\rho b$.  On $0\le s\le1$
write
\[
 \psi(s)=\frac12\min\{s^2,(1-s)^2\},
 \qquad \overline\psi=\int_0^1\psi(s)\,ds=\frac1{24}.
\]
For a good cell $I_j=[\theta_j,\theta_{j+1}]$ define the normalized
bubble
\[
 p_j(\theta_j+h_js)=g^2\psi(s),
 \qquad p_{\rm G}=\sum_{j\in\mathcal G}p_j.
\]
The amplitude error
\[
 f:=q_{\rm G}-p_{\rm G}
 =\sum_{j\in\mathcal G}(h_j^2-g^2)
   \psi\!\left(\frac{\,\cdot-\theta_j\,}{h_j}\right)\mathbf1_{I_j}
\]
is even on every good target cell and is zero on every bad target cell.
It therefore satisfies the self-zero hypothesis \eqref{ZG:eq:selfzero}.
Since $\int_0^1\psi^2=1/320$ and $h_j,g\asymp_\rho b$,
\begin{align}
 \sum_i h_i|\Hil f(c_i)|^2
 &\le C\|f\|_2^2\notag\\
 &=\frac1{320}\sum_{j\in\mathcal G}
 h_j(h_j^2-g^2)^2
 \le C_\rho b^3\sum_{j\in\mathcal G}(h_j-g)^2.
 \label{ZG:eq:amplitude-error}
\end{align}

It remains to estimate $p_{\rm G}$.  Decompose it as
\begin{equation}\label{ZG:eq:normalized-decomposition}
 p_{\rm G}=g^2\overline\psi\,\mathbf1_{G^*}+r_{\rm G},
 \qquad
 r_{\rm G}|_{I_j}
 =g^2\left\{\psi\!\left(\frac{\,\cdot-\theta_j\,}{h_j}\right)
             -\overline\psi\right\}.
\end{equation}
Because $\Hil\mathbf1=0$,
$\Hil\mathbf1_{G^*}=-\Hil\mathbf1_B$.  The function $\mathbf1_B$ is
constant on every target cell, hence also satisfies
\eqref{ZG:eq:selfzero}.  The adaptive sampling estimate gives
\begin{equation}\label{ZG:eq:mean-gap}
 \sum_i h_i|\Hil(g^2\overline\psi\,\mathbf1_{G^*})(c_i)|^2
 \le Cg^4\|\mathbf1_B\|_2^2
 =Cg^4E.
\end{equation}

We now prove the same bound for $r_{\rm G}$ on the bad targets.  Fix a
maximal gap $B_\nu=(a,a+e)$ and $x\in B_\nu$.  Put
\[
 r_-=x-a,\qquad r_+=a+e-x.
\]
The two good cells adjacent to $a$ and $a+e$ have lengths comparable
with $g$.  On either adjacent cell the profile
$\psi-\overline\psi$ has the common endpoint value
$-\overline\psi$; its difference from that endpoint is $\psi$ and
vanishes quadratically.  Since the periodic Hilbert kernel
$K(t)=(2\pi)^{-1}\cot(t/2)$ is odd, pairing the two adjacent cells gives
\begin{equation}\label{ZG:eq:paired-log}
 \left|\int_{I_- \cup I_+}K(x-y)r_{\rm G}(y)\,dy\right|
 \le C_\rho g^2
 \left(1+\left|\log\frac{r_-}{r_+}\right|\right).
\end{equation}
Indeed, the two constant endpoint pieces are
\[
 -g^2\overline\psi
 \left\{\int_0^{h_-}K(r_-+u)\,du
       -\int_0^{h_+}K(r_++u)\,du\right\}.
\]
Lift the gap and its two adjacent cells to the real arc; this is
unambiguous because $E\le\pi/4$ and the fan is sufficiently fine.  Then
\[
 \int_0^hK(t+u)\,du
 =\frac1\pi\log\frac{\sin((t+h)/2)}{\sin(t/2)}
\]
exactly.  On this arc $\sin(v/2)\asymp v$, while
$h_\pm\asymp_\rho g$.  Moreover
$u\mapsto\log(1+e^{-u})$ is $1$-Lipschitz.  Applying this with
$u=\log(t/h)$ shows that the difference of the two logarithms is bounded
by $C_\rho+|\log(r_-/r_+)|$, which proves the constant part of
\eqref{ZG:eq:paired-log}.  For the two remaining pieces use
$\psi(u/h)\le C\min\{(u/h)^2,((h-u)/h)^2\}$; their integrals against
$|K(t+u)|$ are uniformly bounded.

Every other good cell $J$ is at distance at least $c_\rho g$ from $x$.
On $J$ the function $r_{\rm G}$ has zero mass and zero first moment
about the midpoint.  Two Taylor subtractions and
$|K''(t)|\le C\operatorname {dist}(t,2\pi\mathbb Z)^{-3}$ give
\[
 \left|\int_JK(x-y)r_{\rm G}(y)\,dy\right|
 \le \frac{C_\rho g^5}{\operatorname {dist}(x,J)^3}.
\]
The $k$th good cell in either cyclic direction is at distance at least
$c_\rho kg$; summing $k^{-3}$ and combining with
\eqref{ZG:eq:paired-log} yields
\begin{equation}\label{ZG:eq:gap-pointwise}
 |\Hil r_{\rm G}(x)|
 \le C_\rho g^2
 \left(1+\left|\log\frac{r_-}{r_+}\right|\right).
\end{equation}

Finally, for an arbitrary subdivision of $(a,a+e)$ into target cells
with midpoints $c_i$,
\begin{equation}\label{ZG:eq:midpoint-log-square}
 \sum_{I_i\subset B_\nu}h_i
 \left(1+\left|\log
 \frac{c_i-a}{a+e-c_i}\right|\right)^2
 \le Ce.
\end{equation}
To see this, scale to $(0,1)$.  On either half interval the displayed
profile is monotone, so the value at a cell midpoint is bounded by twice
its integral over the half of that cell nearer the corresponding
endpoint.  A cell crossing $1/2$ contributes $O(h_i)$.  The remaining
claim is the elementary finite integral
\[
 \int_0^1\left(1+\left|\log\frac{s}{1-s}\right|\right)^2ds<\infty.
\]
Equations \eqref{ZG:eq:gap-pointwise}--\eqref{ZG:eq:midpoint-log-square},
summed over the disjoint gaps, prove
\begin{equation}\label{ZG:eq:oscillation-gap}
 \sum_{i\in\mathcal B}h_i|\Hil r_{\rm G}(c_i)|^2
 \le C_\rho g^4E.
\end{equation}
Combining \eqref{ZG:eq:amplitude-error}, \eqref{ZG:eq:mean-gap}, and
\eqref{ZG:eq:oscillation-gap}, and using $g\asymp_\rho b$ and
$h_j-g=\widehat h_j-b$, proves \eqref{ZG:eq:gap-target}.
\end{proof}

We can now state the corrected replacement for the stage 155 opening
square.

\begin{lemma}[Zero-gap skeleton square]\label{ZG:lem:skeleton}
Let good source
cells have lengths comparable with $b$, let disjoint
zero-source gaps have total length $E$, and let the target measure put
the actual cell mass at every good or bad midpoint.  If the configuration
is obtained from the collapsed regular $M$-fan by varying the good
lengths to $h_i$ and opening the gaps, then
\begin{equation}\label{ZG:eq:skeleton}
 \boxed{
 \sum_i h_i|\Hil q_{\rm G}(c_i)|^2
 \le C_\rho\left\{
 b^3\sum_{j\in\mathcal G}(h_j-b)^2+b^4E\right\}.}
\end{equation}
The constant is independent of the number of gaps, the number of
target cells inside a gap, and every gap ratio.
\end{lemma}

\begin{proof}
Let $\gamma=E/M$ and $\widehat h_j=h_j+\gamma$.  Then
$\sum_{j\in\mathcal G}\widehat h_j=2\pi$ and
\[
 \sum_{j\in\mathcal G}(\widehat h_j-b)^2
 =\sum_{j\in\mathcal G}(h_j-b)^2-\frac{E^2}{M}
 \le\sum_{j\in\mathcal G}(h_j-b)^2.
\]
The complete fan $\widehat h$ is comparable with $b$.  The variable-grid
$T1$ theorem of stage 154, applied on that complete fan, gives
\[
 \sum_{i\in\mathcal G}h_i
 |\Hil\widehat q(\widehat c_i)|^2
 \le C_\rho b^3
 \sum_{j\in\mathcal G}(h_j-b)^2,
\]
because $h_i\asymp_\rho b$.  This is precisely the square of the
reference vector on the good target atoms.  Lemma~\ref{ZG:lem:opening}
bounds the good-target opening error by $C_\rho b^4E$, and
Lemma~\ref{ZG:lem:gap-target} bounds all bad target atoms.  Since
\[
 \sum_{j\in\mathcal G}(\widehat h_j-b)^2
 \le\sum_{j\in\mathcal G}(h_j-b)^2,
\]
the inequality $|x+y|^2\le2|x|^2+2|y|^2$ proves
\eqref{ZG:eq:skeleton}.
\end{proof}

\subsection{AVS, HQ, and DF}

\begin{theorem}[Adaptive variance sampling]
\label{ZG:thm:AVS}
For every
positive sufficiently fine fan,
\begin{equation}\label{ZG:eq:AVS}
 \boxed{
 T(h)\le C\Dvar(h)\le CS\Jfour(h).}
\end{equation}
\end{theorem}

\begin{proof}
The far-stratum region is Theorem 4.1 of stage 153.  In the near region,
combine \eqref{ZG:eq:Tsplit}, Lemma~\ref{ZG:lem:bad}, and
Lemma~\ref{ZG:lem:skeleton}, and then apply the static budget
\eqref{ZG:eq:budget}.  This gives $T\le C_\rho\Dvar$.  Stage 153 proves
$\Dvar\le S\Jfour$.
\end{proof}

\begin{corollary}[Pure midpoint row DF]\label{ZG:cor:DF}
The pure cubic Bregman row DF holds unconditionally.
\end{corollary}

\begin{proof}
Stage 144 gives the exact identity
\[
 \Cnull(q_h^3)-\Cnull(q_{a\mathbf1}^3)
 =\frac1{720}\left(\sum_i h_i^5-Na^5\right)
  +\frac1{576}\Dvar(h)-T(h),
 \qquad a=\frac{2\pi}{N}.
\]
The linear part of the fifth-power Bregman difference cancels because
$\sum_i(h_i-a)=0$.  Taylor's formula and $h_i,a\le S$ give
\[
 0\le\sum_i h_i^5-Na^5
 \le CS\sum_i(h_i-a)^2(h_i^2+2ah_i+3a^2)
 =CS\Jfour(h).
\]
Also $\Dvar\le S\Jfour$ by stage 153, and
$0\le T\le C\Dvar$ by Theorem~\ref{ZG:thm:AVS}.  Taking absolute values in
the exact identity proves DF with no omitted sign case.
\end{proof}
\endgroup
